\chapter{Effects with a Mata fallback}\label{app:mata}

\nwsaom{} simulates the model in a compiled plugin. The first effects were also written in Mata, and for them the program keeps working when no plugin is available for the platform (or the plugin is too old). The Mata simulator is, however, far too slow for the long simulations that estimation needs: a short two-wave fit that takes well under a second with the plugin takes about 30 seconds in Mata, and the larger models of this manual would take hours. For this reason the effects that were added later exist only in the plugin; writing a second, slower version of each of them would add no capability that a user could actually use. If you ask for such an effect and the plugin cannot be used, \nwsaom{} stops with a message that says so. After a fit, \texttt{e(engine)} tells which simulator ran.

\section{Effects that run in Mata as well}

\begin{description}
\item[Network effects, structural] \opt{outdegree}, \opt{reciprocity}, \opt{transtrip}, \opt{transmedtrip}, \opt{transties}, \opt{transrectrip}, \opt{cycle3}, \opt{cycle4}, \opt{balance}, \opt{between}, \opt{indegpopularity}, \opt{outpopularity}, \opt{outactivity}, \opt{inactivity}, \opt{outactsqrt}, \opt{inpop}, \opt{outpop}, \opt{inact}, \opt{isolatepop}, the shared-partner family \opt{gwesp}, \opt{gwespfb}, \opt{gwespbf}, \opt{gwespbb}, \opt{gwesprr}, \opt{gwdspff}, \opt{gwdspfb}, \opt{nbrdist2}, \opt{dblnbrdist2}, \opt{sqrtcycle4}, the degree-shape effects \opt{outtrunc()}, \opt{outmore()}, \opt{outthreshold()}, \opt{outinv()}, \opt{outsqinv()}, the reciprocal-degree effects \opt{recipop}, \opt{recipopsqrt}, \opt{reciact}, \opt{reciactsqrt}, the assortativity effects \opt{outoutass}, \opt{ininass}, \opt{inoutass}, \opt{outinass}, and the isolate effects \opt{antiiso}, \opt{antiiniso}, \opt{antiiniso2}, \opt{outiso}, \opt{inplus3}, \opt{isolatenet}.
\item[Network effects of covariates] \opt{nodematch()} (\texttt{samex}), \opt{nodecov()}, \opt{nodeicov()}, \opt{nodeocov()} (\texttt{altx}, \texttt{egox}), \opt{simcov()} (\texttt{simx}), \opt{egoxaltx}, \opt{samexrecip}, \opt{simrecipx}, \opt{higher}, \opt{absdiffx}, dyadic covariates \opt{dyadcov()}, \opt{dyadcovrecip()}, and interactions \opt{interact()}.
\item[Endowment and creation] \opt{endow()} and \opt{creation()} of the main network effects.
\item[Behavior effects] \opt{linear}, \opt{quadratic}, \opt{avalt}, \opt{avsim} (with their endowment and creation versions), the first set of \opt{beheff()} tokens (degree and alter effects of the network on the behavior, \opt{totalt}, \opt{maxalt}, \opt{minalt} and the similarity family), and \opt{behsim}, \opt{behegox}, \opt{behaltx}, \opt{behegoplusalt} for the network effects of the behavior itself.
\item[Models and data options] \opt{ratecov()}, \opt{tvcov()}, \opt{present()}, \opt{missnet()}, \opt{missbeh()}, \opt{structural()}, several behaviors in \opt{behavior()}, \opt{waves()}, \opt{unconditional}, and \texttt{estat gof}.
\end{description}

\section{Effects that need the plugin}

The following need the plugin, and \nwsaom{} stops with an error if it cannot be used:
\begin{itemize}
\item the covariate-by-triad, popularity and activity effects \opt{samextranstrip()}, \opt{diffxtranstrip()}, \opt{simxtranstrip()}, \opt{homxtranstrip()}, \opt{samexoutact()}, \opt{diffxoutact()}, \opt{samexreciact()}, \opt{diffxreciact()}, \opt{homxoutact()}, \opt{altxoutact()}, \opt{inpopx()}, \opt{outpopx()}, \opt{inactx()}, \opt{outactx()}, \opt{samexinpop()}, \opt{diffxinpop()};
\item the centred and combined degree effects \opt{inpopc}, \opt{outpopc}, \opt{outactc}, \opt{inactc}, \opt{degplus}, \opt{degplusc}, \opt{avdeg()}, \opt{transtriads};
\item the threshold and contrast effects of covariates \opt{egoplusaltsqx()}, \opt{altsqx()}, \opt{diffsqx()}, \opt{egodiffx()}, \opt{althigheregox()}, \opt{unequalx()}, \opt{egolthresholdx()}, \opt{egorthresholdx()}, \opt{altlthresholdx()}, \opt{altrthresholdx()}, \opt{degabsdiffx()}, \opt{degposdiffx()}, \opt{degnegdiffx()}, \opt{covthreshold()};
\item the structural rate effects \opt{outrate}, \opt{inrate}, \opt{reciprate} and their inverse and logarithmic versions;
\item the option families \opt{netx()}, \opt{ownx()}, \opt{dyadx()}, \opt{tvdyadcov()}, \opt{tvdyadcovrecip()}, \opt{tvdyadx()}, \opt{netbeh()}, \opt{behx()}, \opt{behw()}, \opt{tvbehw()}, \opt{behnet2()}, \opt{behshape()}, \opt{behfrom()}, \opt{behavxalt()}, \opt{behavsamexalt()}, \opt{behtotxalt()}, \opt{behtotsamexalt()} and the later \opt{beheff()} tokens;
\item co-evolution of a non-directed network with a behavior, two-mode networks, and \texttt{nwsaom multinet};
\item \opt{estat importance}, \opt{estat simulate} and \opt{nwsaom simulate} use the plugin for the models they simulate or evaluate that need it.
\end{itemize}
The help file marks the options of the second group with \emph{native plugin only}. The list above describes the situation when the manual was written; a new version of \nwsaom{} adds new effects to the second group only.
